\documentclass[11pt,a4paper]{letter}
\usepackage[T1]{fontenc}
\usepackage[utf8]{inputenc}
\usepackage[left=2.5cm,right=2.5cm,top=2.5cm,bottom=2.5cm]{geometry}
\usepackage{hyperref}

%% --- Zenodo replication-package DOI (single point of update) ---
\newcommand{\zenododoi}{10.5281/zenodo.21939600}
\newcommand{\zenododoiurl}{https://doi.org/\zenododoi}

\signature{Rodrigo Almeida de Oliveira, M.Sc. \\ PhD candidate, PPGIa \\ Pontif\'icia Universidade Cat\'olica do Paran\'a (PUCPR)}
\address{Rodrigo Almeida de Oliveira, M.Sc. \\ Pontif\'icia Universidade Cat\'olica do Paran\'a \\ Programa de P\'os-Gradua\c c\~ao em Inform\'atica (PPGIa) \\ Curitiba, PR, Brazil \\ rodrigo.almeida1@pucpr.edu.br}

\begin{document}

\begin{letter}{Editor-in-Chief \\ \emph{Information and Software Technology} \\ Elsevier}

\opening{Dear Editors,}

On behalf of all co-authors (Rodrigo Almeida de Oliveira, Juliano de Paulo
Ribeiro, and Prof.\ Edson Emilio Scalabrin), I am pleased to submit our
manuscript titled
\textbf{``Process Mining and Stochastic Modeling for Software Process
Forecasting: A Systematic Mapping of Stored Family Codes''}
for consideration in \emph{Information and Software Technology}.

\medskip\noindent\textbf{Novelty and contribution.}
This manuscript reports a systematic mapping study at the
intersection of process mining (PM), stochastic modeling, and machine
learning applied to software process forecasting. The contribution is
threefold:
\begin{enumerate}
\item \textbf{Methodological:} an LLM screened the records (Claude Haiku 4.5
      as primary screener, Claude Sonnet 4.6 as verifier). On the parseable
      cross-model sample, binary Cohen's $\kappa$ is 0.695 at
      title/abstract (420 of 468 papers; 48 verifier responses had no
      category because the abstract was absent) and 0.694 at the second
      pass ($n=177$). That agreement is not a human validation. The
      pre-registered check is a two-rater sample. Recall on the
      working-set sample is 0.630 (29 of 46). It is not a full-text
      Recall: a PDF was read for 0.5\% of the 886 records. The round-1
      FT row (Recall 0.246; $\kappa = 0.122$) is withdrawn. Auxiliary Recall after the
      stratum re-screen is 0.725 (235 of 324). Both are below the
      pre-registered bar of 0.90. The two reviewers agreed with 92\% of
      the suggested labels on the remaining records. The re-screen sheets
      (4{,}005 rows) match the suggestion on every row, so that share is
      not marked there. Rebuilt from the title or abstract, it is 32 of
      the 338 coded includes (9.5\%), and those 32 are removed.
\item \textbf{Empirical:} the corpus of this revision is 456 studies: 306
      that remain after 32 coded studies are removed from the list of 338,
      plus the 150 added studies (ADD-150) that remain after an
      eligibility pass removed 18 of 168. The family codes are reported
      on WS-151, the 151 working-set studies inside the 306, a cut
      disjoint from ADD-150.
      One study
      in the 380 with at least one family mark (PRIMAD) matches process
      mining, a stochastic model, and forecasting. Seventy-six added
      studies, read from the title alone, are outside that fraction.
      F2 and F5 are withdrawn. F1, F3, and F4 are stored codes and are not
      findings. The agenda is not inferred from them.
\item \textbf{Synthetic:} the SPMF taxonomy and a four-direction research
      agenda (RA1--RA4), proposed by the authors and not inferred from
      F1--F5. PATHCAST is future work. It is not a result of
      this mapping.
\end{enumerate}

\medskip\noindent\textbf{Fit with IST scope.}
The manuscript squarely fits IST's scope on empirical software engineering
methods, mining software repositories, and process modeling. It complements
the journal's growing body of evidence-based secondary studies while
contributing to the methodological literature on LLM-assisted reviews.

\medskip\noindent\textbf{Reproducibility.}
A complete replication package is openly available at Zenodo
(\href{\zenododoiurl}{DOI~\zenododoi}),
including all search strings, screening prompts, JSON schemas, decision
CSVs at every pipeline stage, the 456-study included list
(\texttt{included\_studies\_506.csv}), the cross-model $\kappa$ scripts,
and the structured extraction tables.

\medskip\noindent\textbf{Funding.}
This study was financed in part by CAPES (Finance Code 001) through a
doctoral scholarship granted to the corresponding author within the
Graduate Program in Informatics (PPGIa) at PUCPR.

\medskip\noindent\textbf{Prior submission and conflicts.}
The manuscript is original work, has not been published elsewhere, and is
not under consideration by any other journal. All authors declare no
competing financial interests. The use of generative AI tools in the
research is declared in the manuscript per Elsevier policy. AI tools are
\emph{not} listed as authors and bear no responsibility for editorial
judgements.

\medskip\noindent\textbf{Suggested reviewers.}
\begin{itemize}
\item \textbf{Wil M.\ P.\ van der Aalst} --- RWTH Aachen University,
      foundational PM researcher (\texttt{wvdaalst@pads.rwth-aachen.de})
\item \textbf{Maria Teresa Baldassarre} --- University of Bari, empirical
      SE and secondary studies
\item \textbf{Eugene Syriani} --- University of Montreal, LLM-assisted SLR
      methodology
\item \textbf{Tore Dyb{\aa}} --- SINTEF / University of Oslo, evidence-based
      SE
\end{itemize}

\medskip\noindent\textbf{Reviewers to avoid.}
None on basis of co-authorship within the past five years. Specific
exclusions: thesis advisors and committee members at PUCPR.

\closing{Sincerely,}

\end{letter}

\end{document}
